\documentclass[12pt,a4paper]{article}
\usepackage[utf8]{inputenc}
\usepackage[T1]{fontenc}
\usepackage[french]{babel}
\usepackage{geometry}
\usepackage{amsmath}
\usepackage{array}
\usepackage{tabularx}
\usepackage{xcolor}
\usepackage{colortbl}
\usepackage{tikz}
\usepackage{pgfplots}
\usepackage{fancyhdr}
\usepackage{enumitem}
\usepackage{tcolorbox}
\usepackage{amssymb}

\geometry{margin=2cm, top=3cm, bottom=2.5cm}
\setlength{\headheight}{24pt}
\pgfplotsset{compat=1.18}
\tcbuselibrary{skins, breakable}

%--------------------------------------------------
% Couleurs
%--------------------------------------------------
\definecolor{mainblue}{RGB}{0, 102, 153}
\definecolor{lightblue}{RGB}{220, 238, 250}
\definecolor{graylight}{RGB}{245, 245, 245}
\definecolor{grayborder}{RGB}{160, 160, 160}
\definecolor{myorange}{RGB}{200, 90, 0}

%--------------------------------------------------
% En-tête / Pied de page
%--------------------------------------------------
\pagestyle{fancy}
\fancyhf{}
\fancyhead[L]{\small\textcolor{mainblue}{\textbf{Physique-Chimie -- Terminale Bac Pro (Groupement 1,2 et 4)}}}
\fancyhead[R]{\small\textcolor{mainblue}{\textbf{Moteurs \'{e}lectriques}}}
\fancyfoot[C]{\small\thepage}
\renewcommand{\headrulewidth}{1pt}
\renewcommand{\headrule}{\hbox to\headwidth{\color{mainblue}\leaders\hrule height \headrulewidth\hfill}}

%--------------------------------------------------
% Boîtes tcolorbox
%--------------------------------------------------
\newtcolorbox{contexte}{
  colback=lightblue, colframe=mainblue,
  fonttitle=\bfseries\small\color{white},
  title=Contexte,
  arc=4pt, boxrule=1pt,
  colbacktitle=mainblue
}

\newtcolorbox{docbox}[1]{
  colback=graylight, colframe=grayborder,
  fonttitle=\bfseries\small,
  title=#1,
  arc=3pt, boxrule=0.7pt
}

\newtcolorbox{competences}{
  colback=white, colframe=myorange,
  fonttitle=\bfseries\small\color{white},
  title=Comp\'{e}tences \'{e}valu\'{e}es (r\'{e}f\'{e}rentiel BO),
  arc=4pt, boxrule=1pt,
  colbacktitle=myorange
}

\newtcolorbox{essentiel}{
  colback=mainblue!8, colframe=mainblue,
  arc=4pt, boxrule=0.8pt
}

%--------------------------------------------------
% Commandes utilitaires
%--------------------------------------------------
\newcounter{nquestion}
\newcommand{\Q}[1]{%
  \stepcounter{nquestion}%
  \vspace{0.4cm}
  \noindent\textcolor{mainblue}{\textbf{$\blacktriangleright$ Question \thenquestion.}} \textbf{#1}
  \par
}

\newcommand{\partie}[2]{%
  \vspace{0.6cm}
  \noindent
  \begin{tcolorbox}[colback=mainblue, colframe=mainblue, arc=3pt, boxrule=0pt,
    left=6pt, right=6pt, top=4pt, bottom=4pt]
    \textcolor{white}{\textbf{Partie #1 -- #2}}
  \end{tcolorbox}
  \vspace{0.1cm}
}

\newcommand{\lignereponse}{\vspace{0.55cm}\noindent\dotfill\par}

%==================================================
\begin{document}
%==================================================

%--------------------------------------------------
% Titre
%--------------------------------------------------


\begin{center}
  {\Large\bfseries\color{mainblue}
    {Obtenir de l'\'{e}nergie m\'{e}canique \`{a} l'aide d'un moteur \'{e}lectrique} synchrone ou asynchrone}\\[0.4cm]
  \begin{tabular}{@{}lcr@{}}
    \textbf{Dur\'{e}e :} 1\,h
    & \hspace{4cm}
    & \textbf{Classe :} Terminale Bac Pro
  \end{tabular}
\end{center}

{\color{mainblue}\rule{\linewidth}{1pt}}

\vspace{0.3cm}
\textbf{Nom :} \dotfill \hspace{2cm} \textbf{Pr\'{e}nom :} \dotfill \hspace{2cm} \textbf{Classe :} \dotfill

\vspace{0.3cm}
{\color{mainblue}\rule{\linewidth}{1pt}}
\vspace{0.4cm}

%--------------------------------------------------
% Compétences
%--------------------------------------------------
\begin{competences}
\begin{itemize}[noitemsep, leftmargin=*]
  \item Reconna\^{i}tre un moteur \`{a} courant continu et un moteur asynchrone \`{a} partir de sa plaque signal\'{e}tique.
  \item Mettre en \'{e}vidence le principe de conversion d'\'{e}nergie \'{e}lectrom\'{e}canique par un bilan de puissance.
  \item V\'{e}rifier l'influence de la tension d'alimentation sur la fr\'{e}quence de rotation (moteur CC).
  \item V\'{e}rifier l'influence de la fr\'{e}quence de la tension d'alimentation sur la fr\'{e}quence de rotation (moteur asynchrone).
\end{itemize}
\end{competences}

\vspace{0.4cm}

%--------------------------------------------------
% Contexte
%--------------------------------------------------
\begin{contexte}
Dans un atelier de production, deux moteurs \'{e}lectriques entra\^{i}nent diff\'{e}rentes machines~: un
\textbf{moteur M1} pour une perceuse \`{a} colonne \`{a} vitesse variable, et un \textbf{moteur M2} pour
entra\^{i}ner un compresseur d'air. En tant que technicien(ne), vous devez identifier ces moteurs,
analyser leurs caract\'{e}ristiques et comprendre comment r\'{e}gler leur vitesse de rotation.
\end{contexte}

\vspace{0.2cm}

%==================================================
\partie{A}{Identification des moteurs \textnormal{(15 min)}}
%==================================================

\begin{docbox}{Document 1 -- Plaque signal\'{e}tique du moteur M1}
\begin{center}
\renewcommand{\arraystretch}{1.3}
\begin{tabular}{|c|c||c|c|}
\hline
\multicolumn{4}{|c|}{\textbf{MOTEUR M1 -- Leroy-Somer LS80L}} \\
\hline
\textbf{U} & 24\,V \textbf{DC} & \textbf{I} & 4{,}2\,A \\
\hline
\textbf{P} & 90\,W & \textbf{n} & 1500\,tr/min \\
\hline
\textbf{Classe isolation} & F & \textbf{IP} & 54 \\
\hline
\end{tabular}
\end{center}
\end{docbox}

\vspace{0.3cm}

\begin{docbox}{Document 2 -- Plaque signal\'{e}tique du moteur M2}
\begin{center}
\renewcommand{\arraystretch}{1.3}
\begin{tabular}{|c|c||c|c|}
\hline
\multicolumn{4}{|c|}{\textbf{MOTEUR M2 -- Siemens 1LA7 083-2AA}} \\
\hline
\textbf{U} & 230/400\,V \textbf{AC} & \textbf{f} & 50\,Hz \\
\hline
\textbf{I} & 3{,}5\,/\,2{,}0\,A & \textbf{P} & 750\,W \\
\hline
$\mathbf{\cos\,\varphi}$ & 0{,}83 & \textbf{n} & 2840\,tr/min \\
\hline
\textbf{Couplage} & $\triangle$\,/\,Y & \textbf{IP} & 55 \\
\hline
\end{tabular}
\end{center}
\end{docbox}

\vspace{0.3cm}

\begin{docbox}{Document 3 -- Rappel : les deux cat\'{e}gories de moteurs \'{e}lectriques}
Un moteur \'{e}lectrique est un \textbf{convertisseur \'{e}lectrom\'{e}canique} : il transforme l'\'{e}nergie
\'{e}lectrique en \'{e}nergie m\'{e}canique.

\medskip
\begin{tabular}{@{}p{7.2cm}@{\hspace{0.6cm}}p{7.2cm}@{}}
\textbf{Moteur \`{a} courant continu (MCC)} &
\textbf{Moteur asynchrone (triphas\'{e})} \\[0.2cm]
Alimentation : tension \emph{continue} (DC).
La vitesse varie avec la tension d'alimentation.
La plaque indique une tension en \textbf{V DC}.
&
Alimentation : tension \emph{alternative} (AC), souvent triphas\'{e}e.
La vitesse d\'{e}pend de la fr\'{e}quence du r\'{e}seau.
La plaque indique une fr\'{e}quence en \textbf{Hz}.
\end{tabular}
\end{docbox}

\Q{\`{A} partir des documents 1, 2 et 3, identifiez la cat\'{e}gorie de chaque moteur
(courant continu ou asynchrone). Justifiez en relevant les informations pertinentes sur
chaque plaque signal\'{e}tique.}

\lignereponse\lignereponse\lignereponse\lignereponse

\Q{Pour le moteur M2, la vitesse de synchronisme est donn\'{e}e par $n_s = \dfrac{60\,f}{p}$
o\`{u} $p$ est le nombre de paires de p\^{o}les. Sachant que $n_s = 3000$\,tr/min pour $f = 50$\,Hz,
calculez le glissement $g = \dfrac{n_s - n}{n_s}$ et exprimez-le en \%.}

\lignereponse\lignereponse\lignereponse

\newpage

%==================================================
\partie{B}{Bilan de puissance et rendement \textnormal{(20 min)}}
%==================================================

\begin{docbox}{Document 4 -- Sch\'{e}ma de conversion \'{e}lectrom\'{e}canique}
\begin{center}
\begin{tikzpicture}[>=stealth, scale=0.95]
  % ---- Bloc central : le moteur ----
  \draw[thick, fill=lightblue, rounded corners=6pt] (2.5,0) rectangle (8.5,2);
  \node[font=\bfseries\large] at (5.5,1.4) {MOTEUR};
  \node[font=\small] at (5.5,0.7) {Convertisseur \'{e}lectrom\'{e}canique};

  % ---- Flèche entrée : puissance électrique (gauche → moteur) ----
  \draw[->, line width=2pt, color=mainblue] (-1,1) -- (2.5,1);
  \node[above, color=mainblue, font=\bfseries] at (0.85,1.05) {$P_{\text{\'{e}lec}}$};
  \node[below, color=mainblue, font=\small] at (0.85,0.85) {\'{e}nergie \'{e}lectrique};

  % ---- Flèche sortie : puissance mécanique (moteur → droite) ----
  \draw[->, line width=2pt, color=myorange] (8.5,1) -- (12.3,1);
  \node[above, color=myorange, font=\bfseries] at (10.35,1.05) {$P_{\text{m\'{e}c}}$};
  \node[below, color=myorange, font=\small] at (10.35,0.85) {\'{e}nergie m\'{e}canique};

  % ---- Flèche pertes : chaleur dissipée (moteur → haut) ----
  \draw[->, line width=1.8pt, color=red!70] (5.5,2) -- (5.5,3.3);
  \node[right, color=red!70, font=\small] at (5.6,2.65) {$P_{\text{pertes}}$ (chaleur)};
\end{tikzpicture}
\end{center}
\vspace{0.15cm}
\[
  P_{\text{\'{e}lec}} = P_{\text{m\'{e}c}} + P_{\text{pertes}}
  \qquad\qquad
  \eta = \frac{P_{\text{m\'{e}c}}}{P_{\text{\'{e}lec}}}
\]
\begin{itemize}[noitemsep, leftmargin=*]
  \item $P_{\text{\'{e}lec}} = U \times I$ \quad (moteur CC, unit\'{e} : W)
  \item $P_{\text{m\'{e}c}} = C \times \omega$ \quad avec $\omega = \dfrac{2\pi\,n}{60}$ (rad/s), $C$ : couple utile (N\,.\,m)
  \item $\eta$ : rendement, sans unit\'{e}, compris entre 0 et 1
\end{itemize}
\end{docbox}

\vspace{0.3cm}

\begin{docbox}{Document 5 -- Mesures effectu\'{e}es sur le moteur M1 en charge}
\begin{center}
\renewcommand{\arraystretch}{1.35}
\begin{tabular}{|l|c|}
\hline
Tension d'alimentation & $U = 24$\,V \\
\hline
Intensit\'{e} du courant & $I = 3{,}8$\,A \\
\hline
Vitesse de rotation mesur\'{e}e & $n = 1\,420$\,tr/min \\
\hline
Couple utile & $C = 0{,}55$\,N\,.\,m \\
\hline
\end{tabular}
\end{center}
\end{docbox}

\Q{Calculez la puissance \'{e}lectrique absorb\'{e}e $P_{\text{\'{e}lec}}$ par le moteur M1.}

\lignereponse\lignereponse

\Q{Calculez la vitesse angulaire $\omega$ puis la puissance m\'{e}canique utile $P_{\text{m\'{e}c}}$.}

\lignereponse\lignereponse\lignereponse

\newpage

\Q{D\'{e}duisez les pertes $P_{\text{pertes}}$ et le rendement $\eta$ du moteur M1.
Comparez $\eta$ \`{a} la valeur que l'on peut estimer \`{a} partir de la plaque signal\'{e}tique.
Commentez.}

\lignereponse\lignereponse\lignereponse

%==================================================
\partie{C}{Influence des param\`{e}tres \'{e}lectriques sur la vitesse \textnormal{(25 min)}}
%==================================================

\begin{docbox}{Document 6 -- Moteur M1 (CC) : influence de la tension sur la vitesse}
On fait varier la tension d'alimentation $U$ du moteur M1 \`{a} vide et on mesure la fr\'{e}quence de rotation $n$.
\begin{center}
\renewcommand{\arraystretch}{1.3}
\begin{tabular}{|c|c|c|c|c|c|c|}
\hline
\textbf{Tension $U$ (V)} & 4 & 8 & 12 & 16 & 20 & 24 \\
\hline
\textbf{Vitesse $n$ (tr/min)} & 220 & 450 & 680 & 920 & 1\,170 & 1\,430 \\
\hline
\end{tabular}
\end{center}
\end{docbox}

\vspace{0.3cm}

\begin{docbox}{Document 7 -- Moteur M2 (asynchrone) : influence de la fr\'{e}quence sur la vitesse}
Un \textbf{variateur de fr\'{e}quence} fait varier la fr\'{e}quence $f$ de la tension d'alimentation
du moteur M2~; on mesure la fr\'{e}quence de rotation $n$.
\begin{center}
\renewcommand{\arraystretch}{1.3}
\begin{tabular}{|c|c|c|c|c|c|}
\hline
\textbf{Fr\'{e}quence $f$ (Hz)} & 10 & 20 & 30 & 40 & 50 \\
\hline
\textbf{Vitesse $n$ (tr/min)} & 560 & 1\,140 & 1\,710 & 2\,280 & 2\,840 \\
\hline
\end{tabular}
\end{center}
\end{docbox}

\vspace{0.3cm}

\Q{Tracez le graphique $n = f(U)$ \`{a} partir du document 6.
D\'{e}crivez l'\'{e}volution de la vitesse en fonction de la tension.}

\begin{center}
\begin{tikzpicture}
\begin{axis}[
  width=13cm, height=6.5cm,
  xlabel={Tension $U$ (V)},
  ylabel={Vitesse $n$ (tr/min)},
  xmin=0, xmax=26,
  ymin=0, ymax=1600,
  grid=both,
  grid style={line width=0.3pt, draw=gray!30},
  major grid style={line width=0.6pt, draw=gray!55},
  xtick={0,2,...,26},
  ytick={0,200,...,1600},
  tick label style={font=\small},
  label style={font=\small\bfseries},
  axis line style={thick},
]
\end{axis}
\end{tikzpicture}
\end{center}

\lignereponse\lignereponse

\Q{Quelle est la cons\'{e}quence pratique pour le contr\^{o}le de la vitesse de la perceuse
entra\^{i}n\'{e}e par le moteur M1 ?}

\lignereponse\lignereponse\lignereponse

\Q{Tracez le graphique $n = f(f_{\text{r\'{e}seau}})$ \`{a} partir du document 7.
D\'{e}crivez l'\'{e}volution de la vitesse en fonction de la fr\'{e}quence.}

\begin{center}
\begin{tikzpicture}
\begin{axis}[
  width=13cm, height=6.5cm,
  xlabel={Fr\'{e}quence du r\'{e}seau $f$ (Hz)},
  ylabel={Vitesse $n$ (tr/min)},
  xmin=0, xmax=55,
  ymin=0, ymax=3200,
  grid=both,
  grid style={line width=0.3pt, draw=gray!30},
  major grid style={line width=0.6pt, draw=gray!55},
  xtick={0,5,...,55},
  ytick={0,400,...,3200},
  tick label style={font=\small},
  label style={font=\small\bfseries},
  axis line style={thick},
]
\end{axis}
\end{tikzpicture}
\end{center}

\lignereponse\lignereponse

\Q{Expliquez le principe du variateur de fr\'{e}quence pour r\'{e}gler la vitesse du moteur
asynchrone M2. Donnez un exemple d'application industrielle.}

\lignereponse\lignereponse\lignereponse

%

\end{document}
